\section{Асинхронна архитектура на изпълнението}
\label{sec:async_execution_architecture}

След като архитектурният модел по време на компилация е дефиниран, следващият въпрос е как библиотеката реализира асинхронното изпълнение като втора основна архитектурна ос. Настоящата глава разглежда основните градивни елементи на асинхронния слой и ги проследява от \code{Task<T>} и композицията, базирана на \code{co\_await}, до \code{thread\_pool}, \code{async\_db<DB>}, \code{IoContext}, съобразеното с корутини поддържане на пул от връзки и интеграцията със специфични за целевата система пътища на изпълнение.

Тук фокусът вече не е върху междинното представяне на заявките (междинно представяне на заявките) като структура по време на компилация, а върху реалната механика на спиране/възобновяване (спиране/възобновяване) на изпълнението. Как синхронният програмен интерфейс на конектора се повдига до асинхронен интерфейс? Къде изнасянето (изнасяне) към набор от нишки е правилната стратегия и къде платформено-специфичният неблокиращ интерфейс е по-точният модел? Как асинхронният слой гарантира коректна употреба на връзките и транзакциите без връщане към интерфейс, управляван от функции за обратно извикване (управляван от функции за обратно извикване)? Отговорът на тези въпроси показва доколко библиотеката е завършена не само като типова, но и като изпълнителна система.

\subsection{Архитектурна позиция на асинхронния слой}

Асинхронният слой е разположен над вече съществуващия синхронен модел \code{db<DB>} и \code{connector\_trait<DB>}. Това е важен архитектурен избор. Вместо да дублира конструктора на заявки или да въвежда втори набор от типове заявки, асинхронната подсистема приема конструираното по време на компилация междинно представяне като готов вход и се занимава единствено с това \emph{как} да бъде изпълнено. По този начин първата архитектурна ос --- типово описание на данните и операциите --- остава отделена от втората ос --- управление на асинхронния жизнен цикъл.

От тази гледна точка асинхронната архитектура има три основни отговорности. Първо, да осигури корутинен носител на състояние на изпълнение и продължение. Второ, да даде универсален механизъм за повдигане на блокиращи конектори към интерфейс, базиран на \code{Task}. Трето, да позволи платформено-специфични неблокиращи свързващи компоненти да участват в същия потребителски синтаксис, без библиотеката да симулира фалшива еднородност там, където драйверите реално са различни.

\subsection{\code{Task<T>} като носител на жизнения цикъл на корутината}

Основната потребителска абстракция е \code{Task<T>}. Нейната роля не се изчерпва с това да „върне бъдещ резултат“. Тя е контейнер за корутинната рамка, обекта-обещание, продължението и политиката за завършване. Реализацията използва \code{initial\_suspend()} и \code{final\_suspend()} така, че корутината да бъде отложена (отложена) по начало и да предава управлението обратно към продължението при завършване. Следователно \code{Task<T>} е едновременно манипулатор на изпълнението и формален договор за това как резултат, изключение и завършване се предават между корутинни контексти.

От инженерна гледна точка са важни три свойства. Първо, \code{operator co\_await()} свързва текущото продължение към изчакваната (изчаквана) задача и позволява естествено вложено композиране чрез \code{co\_await}. Второ, \code{sync\_wait()} осигурява мост към не-корутинен код, което е необходимо както за модулните (модулни) тестове, така и за гранични сценарии, в които приложението още не работи изцяло по корутинен начин (ориентиран към корутини). Трето, \code{detach()} и \code{start\_detached()} позволяват контролирано поведение "изстреляй и забрави" (изстреляй и забрави), при което корутинната рамка се самоунищожава след \code{final\_suspend()}.

Тази конструкция не е декоративна. Тя определя кога дадена асинхронна операция е отложена, кога започва да се изпълнява, кой носи отговорност за жизнения цикъл на рамката и как се разпространяват изключенията. Именно поради това \code{Task<T>} стои в центъра на всички останали асинхронни компоненти: \code{run\_on\_pool()}, \code{async\_db<DB>}, \code{IoContext} изчакваемите обекти и помощниците за транзакции работят през един и същ корутинен договор.

\subsection{\code{thread\_pool} и \code{run\_on\_pool()}: универсалният път на изпълнение}

Универсалната асинхронна стратегия е изнасяне към набор от нишки (изнасяне към набор от нишки). \code{thread\_pool} поддържа ограничен набор от работни нишки и опашка от задачи. Върху него стъпва \code{run\_on\_pool()}, който приема извикваем обект (извикваем обект), публикува го към опашката, спира текущата корутина и я възобновява, когато извикваемият обект завърши. По този начин синхронният програмен интерфейс на конектора може да бъде използван през \code{co\_await}, без потребителският код да се връща към ръчно свързване на обещания/бъдеща стойност (ръчно свързване на обещания/бъдещи стойности).

Критичният детайл е, че \code{run\_on\_pool()} не е просто помощник за „пусни нещо на друга нишка“. Неговият изчакващ обект (изчакващ обект) пази продължението, евентуалното изключение и резултата от извикваемия обект. След изпълнение върху работната нишка продължението се възобновява (възобновяване), а \code{await\_resume()} връща резултат или повдига грешка. Така пътят с изнасяне остава подчинен на същите корутинни правила като платформено-специфичния асинхронен път.

Този универсален модел е особено важен за свързващи компоненти, които нямат естествен платформено-специфичен неблокиращ клиентски интерфейс или при които библиотеката съзнателно не претендира такъв. В тези случаи честният архитектурен избор е блокиращата операция да бъде изолирана в \code{thread\_pool}, а потребителят да получи единен интерфейс, базиран на \code{Task}. Точно това поведение се валидира от \code{RunOnPoolTest.*} и от \code{AsyncDbTest.AsyncSelectRunsOnPoolThread}, където се доказва, че реалното изпълнение се премества извън главната нишка.

\subsection{\code{async\_db<DB>} и диспечиране, ограничено от възможностите}

Най-важната външна асинхронна фасада е \code{async\_db<DB>}. Тя обгръща \code{db<DB>} и запазва познатия синтаксис на заявките, но връща \code{Task<result<...>>}. Архитектурното значение на тази обвивка е, че тя не решава динамично какво да направи по време на изпълнение, а използва ограничаване на база възможности (ограничаване на възможностите) по време на компилация. При \code{operator<<} има \code{if constexpr} разклонение върху \code{has\_capability<DB, cap::supports\_async>}. Ако конекторът декларира платформено-специфична асинхронна възможност, се извиква \code{connector\_trait<DB>::async\_execute()}. В противен случай синхронният път на \code{execute()} се обгръща в \code{run\_on\_pool()}.

\begin{lstlisting}[language=C++,caption={Централното разклонение, ограничено от възможностите, в \\code{async\_db<DB>}},label={lst:async_db_dispatch}]
template <typename Заявка>
[[nodiscard]] auto operator<<(Заявка q)
    -> Task<decltype(std::declval<db<DB>>() << std::declval<Заявка>())>
{
    using ResultType = decltype(std::declval<db<DB>>() << std::declval<Заявка>());

    if constexpr (has_възможност<DB, cap::supports_асинхронен>)
    {
        co_return co_await конектор_trait<DB>::асинхронен_execute(
            db_.connection(), std::move(q));
    }
    else
    {
        co_return co_await run_on_pool(
            *pool_, [this, q = std::move(q)]() mutable -> ResultType {
                return db_ << std::move(q);
            });
    }
}
\end{lstlisting}

Listing~\ref{lst:async_db_dispatch} е централен за цялата асинхронна архитектура. Той показва, че библиотеката не използва регистър по време на изпълнение или базирано на низове диспечиране (базирано на низове диспечиране), а проверка по време на компилация върху възможностите (възможности) на конектора. По този начин потребителският програмен интерфейс остава еднороден, докато вътрешният път на изпълнение остава верен на реалните възможности на дадената целева система.

Помощният метод \code{async\_execute(Query, Params...)} изпълнява сходна свързваща роля за императивни сценарии с явни параметри, а \code{sync\_handle()} позволява контролиран изход към синхронния интерфейс, когато такава граница е необходима. Следователно \code{async\_db<DB>} не е паралелна библиотека, а адаптационен слой между ядрото за ОРС по време на компилация и корутинния модел на изпълнение.

\begin{figure}[H]
\centering
\begin{tikzpicture}[node distance=1.4cm and 1.6cm]
    \node[ormaccent, text width=3.6cm, align=center] (call) {Потребителска корутина\\\code{co\_await (adb << query)}};
    \node[ormblock, below=of call, text width=3.2cm, align=center] (adb) {\code{async\_db<DB>}};
    \node[ormblock, below=of adb, text width=4.3cm, align=center] (gate) {Разклонение по време на компилация\\\code{has\_capability<DB, cap::supports\_async>}};

    \node[ormblock, below left=of gate, text width=3.0cm, align=center] (poolawait) {\code{run\_on\_pool()}};
    \node[ormblock, below=of poolawait, text width=3.4cm, align=center] (синхроненpath) {\code{thread\_pool} +\\\code{db<DB>::execute}};

    \node[ormblock, below right=of gate, text width=3.3cm, align=center] (платформено-специфичен) {\code{connector\_trait<DB>::async\_execute}};
    \node[ormblock, below=of платформено-специфичен, text width=3.6cm, align=center] (iopath) {\code{IoContext} + платформено-специфичен\\автомат с крайни състояния на драйвера};

    \node[ormaccent, below=of $(синхроненpath)!0.5!(iopath)$, text width=3.3cm, align=center] (result) {\code{Task<result<...>>}\\продължение и резултат};

    \draw[ormline] (call) -- (adb);
    \draw[ormline] (adb) -- (gate);
    \draw[ormline] (gate) -- (poolawait);
    \draw[ormline] (poolawait) -- (синхроненpath);
    \draw[ormline] (gate) -- (платформено-специфичен);
    \draw[ormline] (платформено-специфичен) -- (iopath);
    \draw[ormline] (синхроненpath) -- (result);
    \draw[ormline] (iopath) -- (result);
\end{tikzpicture}
\caption{Двата пътя на изпълнение на \code{async\_db<DB>}: универсално изнасяне (изнасяне) и платформено-специфична асинхронна интеграция}
\label{fig:async_dispatch_paths}
\end{figure}

Фигура~\ref{fig:async_dispatch_paths} показва най-важното архитектурно решение в асинхронния слой. Външният интерфейс остава единен, но вътрешният път на изпълнение е ограничен от възможностите (ограничен от възможностите). Това е причината библиотеката едновременно да предлага удобен корутинен програмен интерфейс и да избягва фалшивото твърдение, че всички драйвери са естествено неблокиращи.

\subsection{\code{IoContext} и платформено-специфична неблокираща интеграция}

За драйвери с истински платформено-специфичен асинхронен интерфейс библиотеката не се ограничава до изнасяне (изнасяне) към набор от нишки. Тук влиза в действие \code{IoContext} --- абстракция за цикъл за събития, която предоставя \code{run()}, \code{run\_one()}, \code{stop()}, \code{post()}, \code{schedule()} и най-важното \code{watch\_readable(fd)} и \code{watch\_writable(fd)}. Тези изчакваеми обекти не извършват входно-изходна операция сами по себе си. Те само спират корутината и я възобновяват, когато файловият дескриптор стане готов за четене или запис.

Точно това разделение между готовност и самото I/O действие е решаващо. \code{IoContext} не се опитва да „скрие“ драйвера, а предоставя съобразен с корутините реакторен (реакторен) интерфейс, върху който платформено-специфичната клиентска библиотека може да стъпи. В реализацията на \code{AsyncPostgreSQLDB} това се вижда особено ясно. Асинхронното свързване използва \code{PQconnectPoll()}, а корутината последователно изчаква (\code{await}) \code{watch\_writable()} или \code{watch\_readable()} според състоянието на автомата с крайни състояния на \code{libpq}. При изпълнение на заявка \code{PQsendQueryParams()}, \code{PQflush()}, \code{PQisBusy()} и \code{PQconsumeInput()} се комбинират със същия механизъм от \code{IoContext}.

Следователно платформено-специфичният асинхронен път не е „по-бърз набор от нишки“, а качествено различен изпълнителен модел. Работната нишка не блокира в очакване на мрежова готовност; корутината се спира (спиране), а възобновяването се управлява от наблюдението на цикъла за събития. Това е точно архитектурният модел, който следващата глава развива по свързващи компоненти: PostgreSQL, MySQL, Redis и Cassandra могат да предложат платформено-специфична асинхронна интеграция с различни експлоатационни форми, докато други драйвери остават честно описани като блокиращи и следователно се обслужват от универсалния път за изнасяне.

\subsection{Поддържане на пул от връзки, съобразено с корутини, и помощници за транзакции}

Асинхронната архитектура не приключва с единичното изпълнение на заявка. Реалното приложение трябва да управлява връзки и транзакции при конкурентен достъп. Именно затова библиотеката предоставя \code{async\_connection\_pool<DB, N>}, \code{async\_connection\_guard<DB>}, \code{async\_transaction\_guard<DB>} и фабричната корутина \code{async\_begin\_transaction()}.

\code{async\_connection\_pool<DB, N>} е съобразено с корутини продължение на синхронния слой за поддържане на пул. Неговият метод \code{acquire()} връща \code{Task<async\_connection\_guard<DB>>}. Вътрешно той използва \code{run\_on\_pool()} и условна променлива (условна променлива), за да изчака слот за връзка със състояние \code{Idle}, без извикващата корутина да бъде изложена на блокиращ интерфейс. Полученият предпазител (предпазител) следва идиомата за придобиване на ресурси чрез инициализация (RAII): при разрушаване той връща състоянието на връзката към \code{Idle} и уведомява чакащите.

От своя страна \code{async\_connection\_guard<DB>::get()} не връща сурова връзка, а нов манипулатор \code{async\_db<DB>}. Това е важен архитектурен детайл. Дори след придобиване на конкретен ресурс за връзка потребителят остава в същия корутинен модел на изпълнение и не пада обратно към програмен интерфейс на ниско ниво.

Сходна логика се вижда и при транзакциите. \code{async\_begin\_transaction()} първо изнася (изнасяне-ва) \code{BEGIN} към пула, а след това връща \code{async\_transaction\_guard<DB>}. Методите \code{commit()} и \code{rollback()} са \code{Task<void>} операции, а деструкторът на предпазителя изпълнява защитен \code{ROLLBACK}, ако транзакцията не е финализирана изрично. Този дизайн е едновременно удобен и консервативен: пътят на успеха е съобразен с корутини, а пътят на изчистването (път на изчистването) остава детерминистичен.

Точно тези свойства се валидират от \code{AsyncConnectionPoolTest.ConcurrentAcquires}, \code{AsyncTransactionTest.BeginAndCommit}, \code{AsyncTransactionTest.RollbackOnDestruction} и \code{AsyncTransactionTest.ExplicitRollback}. Следователно помощниците за пул и транзакции не са периферни удобства, а важна част от цялостната архитектура на изпълнение.

\begin{table}[H]
\centering
\small
\caption{Основни асинхронни подсистеми и тяхната кодова и тестова опора}
\label{tab:async_subsystem_evidence}
\begin{tabularx}{\textwidth}{L{2.9cm}L{4.3cm}Z}
\toprule
\textbf{Подсистема} & \textbf{Архитектурна роля} & \textbf{Основни файлове и тестове} \\
\midrule
\code{Task<T>} & жизнен цикъл на корутината, продължение, \code{sync\_wait()}, семантика на отделяне (detach) & \path{lib/include/ORM/асинхронен/task.hpp}, \path{tests/модулни/test_task.cpp} \\
\code{thread\_pool} и \code{run\_on\_pool()} & универсален път за изнасяне (изнасяне) за блокиращи операции & \path{lib/include/ORM/асинхронен/нишка_pool.hpp}, \path{tests/модулни/test_нишка_pool.cpp} \\
\code{async\_db<DB>} & ограничено от възможности диспечиране над \code{db<DB>} & \path{lib/include/ORM/конектор/асинхронен_db.hpp}, \path{tests/модулни/test_асинхронен_конектор.cpp} \\
\code{IoContext} и платформено-специфични асинхронни конектори & готовност на файлови дескриптори в реакторен стил и платформено-специфична неблокираща интеграция & \path{lib/include/ORM/асинхронен/io_context.hpp}, \path{lib/include/ORM/db/конекторs/PostgreSQLDB/postgresql_асинхронен.hpp}, \path{tests/модулни/test_io_context.cpp} \\
Пул и транзакции & съобразена с корутини собственост върху връзките и транзакционна безопасност & \path{lib/include/ORM/db/конекторs/НишкаSafety/асинхронен_нишка_safety.hpp}, \path{tests/модулни/test_асинхронен_конектор.cpp} \\
\bottomrule
\end{tabularx}
\end{table}

\subsection{Архитектурен извод за асинхронния слой}

Асинхронната архитектура на библиотеката не е просто добавена корутинна обвивка около синхронното ОРС ядро. Тя е самостоятелен слой с ясно разграничени отговорности: \code{Task<T>} моделира жизнения цикъл на корутината, \code{run\_on\_pool()} осигурява универсален механизъм за изнасяне (изнасяне), \code{async\_db<DB>} реализира диспечиране, ограничено от възможностите, \code{IoContext} предоставя платформено-специфична интеграция на готовността, а помощниците за пул и транзакции разширяват модела към практически приложими многозадачни сценарии.

Следователно втората архитектурна ос на дипломната работа не се свежда до това „да има асинхронен програмен интерфейс“. Тя показва как формализираният модел за достъп по време на компилация може да бъде изпълняван през два честно разграничени асинхронни пътя на изпълнение --- универсален и платформено-специфичен --- без библиотеката да губи типова дисциплина, безопасност на жизнения цикъл или прозрачност на конектора. Именно този слой прави възможен следващия анализ на изпълнителните модели на свързващите компоненти.
